\documentclass[10pt,a4paper]{amsart}
\usepackage[margin=19mm]{geometry}
\usepackage{amsmath,amssymb,mathtools}
\usepackage[scr=rsfs,cal=euler]{mathalfa}
\usepackage{xfrac}
\usepackage[unicode,hidelinks,bookmarks=false]{hyperref}
\newcommand{\ie}{i.e.\ }
\newcommand{\cf}{cf.\ }
\newcommand{\resp}{respectively\ }
\newcommand{\Subsection}{\subsection}
\providecommand{\nobreakdash}{\nobreak\hbox{-}}
\setlength{\emergencystretch}{2em}
\multlinegap0pt
\usepackage[shortlabels]{enumitem}
\usepackage{xfrac}
\usepackage{nicefrac}
\usepackage{amssymb}
\usepackage{dsfont}
\usepackage{multirow}
\usepackage{algorithm}\usepackage{algpseudocode}
\usepackage{booktabs}
\usepackage{mathtools}
\newtheorem{proposition}{Proposition}
\newtheorem{lemma}{Lemma}
\newtheorem{definition}{Definition}
\newtheorem{theorem}{Theorem}
\newcommand{\CSL}{C^{\,\sf SL}_{q,p,k}}
\newcommand{\Mk}{\bbM^k_\varepsilon} 
\newcommand{\Omegat}{\widetilde{\Omega}} 
\newcommand{\bbE}{\mathbb{E}}
\newcommand{\bbEt}{\widetilde{\bbE}} 
\newcommand{\bbM}{\mathbb{M}}
\newcommand{\bbN}{\mathbb{N}}
\newcommand{\bbP}{\mathbb{P}}
\newcommand{\bbPt}{\widetilde{\bbP}} 
\newcommand{\bbR}{\mathbb{R}}
\newcommand{\cA}{\mathcal{A}}
\newcommand{\cAt}{\widetilde{\cA}} 
\newcommand{\dual}[1]{{#1}^{\prime}}
\newcommand{\eps}{\varepsilon}
\newcommand{\from}{\colon}
\newcommand{\norm}[2]{     \| #1       \|_{ #2 }}
\newcommand{\omegat}{\widetilde{\omega}} 
\newcommand{\rd}{\mathrm{d}}
\newcommand{\rvec}{\boldsymbol{r}}
\newcommand{\tenvec}{U}
\newcommand{\xivec}{\boldsymbol{\xi}} 
\newcommand{\xk}{\otimes^k} 
\newcommand{\xkseps}{\otimes^{k,s}_{\varepsilon_s}}
\title{Monte Carlo convergence rates for $k$th moments in Banach spaces}
\author{Kristin Kirchner, Christoph Schwab}
\date{Source: arXiv:2212.03797v2 (CC BY 4.0)}
\begin{document}
\maketitle

\noindent\emph{Source and licence.} Monte Carlo convergence rates for $k$th moments in Banach spaces. Version arXiv:2212.03797v2 (CC BY 4.0), distributed by arXiv under the Creative Commons Attribution 4.0 International licence, http://creativecommons.org/licenses/by/4.0/, as recorded in the arXiv licence metadata for this exact version and shown on the abstract page https://arxiv.org/abs/2212.03797v2. Full licence notice: \texttt{licenses/arxiv-2212.03797-CC-BY-4.0.txt}.\par\smallskip

\noindent\textit{Editorial scope.} Counted unit: the article's theorem thm:MC-k (source0.tex lines 2084--2120). The article's complete original proof of it is delivered with it (source0.tex lines 2122--2454, one \texttt{proof} environment of 333 lines). No mathematics is written, repaired or re-derived: the statement and the proof are the article's own lines, in the article's own notation, and no retained line is substituted or re-flowed.  Retained uncounted as the source-local context the proof needs: lines 885--896; lines 1397--1405; lines 1409--1473; lines 1666--1689; lines 1691--1713; lines 1755--1791; lines 1846--1860.  Nothing was omitted from any retained span, so the package carries no omission record.  No retained line contains a citation, so the wrapper carries no bibliography and no bibliography entry.  No retained line contains a figure environment, an image-inclusion command or drawing code, so no image is delivered and the package declares no asset dependency.  Editorial formatting only, in addition: an AMS \texttt{amsart} preamble that restates the article's own declarations and macros used by the retained lines, namely the environments \texttt{proposition}, \texttt{lemma}, \texttt{definition}, \texttt{theorem} on separate counters, together with the standard packages those lines need.  The retained environments print in this document as Proposition 1, Lemma 1, Definition 1, Theorem 1 in order of appearance, whereas the article prints the same items with the numbering of its own full document.  The complete native source of the article as distributed in the arXiv e-print for this exact version is bundled unchanged as \texttt{sources/134602/source0.tex}.
\begin{equation}\label{eq:sepsnorm} 
\norm{\tenvec}{\eps_s} 
:= 
\sup\Biggl\{
\biggl| 
\sum_{j=1}^M 
\lambda_j f(x_j)^k
\biggr| 
\; \Bigg| \; 
f\in B_{\dual{E}} 
\Biggr\}, 
\end{equation} 

	\begin{equation}\label{eq:C-rade-gaussian} 
	C_z 
	:= 
	\begin{cases} 
	1 &\text{ if $z$ is a Rademacher family} , 
	\\
	2 &\text{ if $z$ is an orthogaussian family}. 
	\end{cases} 
	\end{equation}

\begin{proposition}\label{prop:rade-gaussian}  
	Let $M\in\bbN$, 
	$x_1,\ldots,x_M\in E$, 
	and $(z_j)_{j=1}^M$ be 
	a Rademacher or 
	orthogaussian family 
	on a complete probability space 
	$(\Omegat,\cAt,\bbPt)$ 
	(with expectation $\bbEt$). 
	% 
	In addition, assume that 
	$G\from [0,\infty) \to [0,\infty)$ 
	is a convex and increasing function. 
	%  
	\begin{enumerate}[leftmargin=1cm] 
		\item\label{prop:rade-gaussian-poly}  
		For integers $k_1, \ldots, k_M\in\bbN$,  
		we have the relation  
		% 
		\begin{equation}\label{eq:prop:rade-gaussian-poly} 
		\bbEt 
		G\Biggl( 
		\sup_{f\in B_{\dual{E}}}  
		\biggl| 
		\sum_{j=1}^M 
		z_j f(x_j)^{k_j} \biggr| 
		\Biggr) 
		\leq 
		\bbEt 
		G\Biggl(  
		2 C_z \, 
		\biggl\| 
		\sum_{j=1}^M 
		z_j k_j \norm{x_j}{E}^{k_j-1} x_j \biggr\|_E 
		\Biggr).  
		\end{equation}
		% 
		\item\label{prop:rade-gaussian-q}   
		For general exponents 
		$q_1,\ldots,q_M\in[1,\infty)$,  
		the following holds: 
		%
		\begin{equation}\label{eq:prop:rade-gaussian-q} 
		\bbEt 
		G\Biggl(  
		\sup_{f\in B_{\dual{E}}}  
		\biggl| 
		\sum_{j=1}^M 
		z_j |f(x_j)|^{q_j} \biggr| 
		\Biggr) 
		\leq 
		\bbEt 
		G\Biggl(  
		2 C_z \, 
		\biggl\| 
		\sum_{j=1}^M 
		z_j q_j \norm{x_j}{E}^{q_j-1}  x_j \biggr\|_E 
		\Biggr). 
		\end{equation} 
		%
	\end{enumerate}
	% 
	Here, the constant $C_z\in\{1,2\}$ is defined 
	as in \eqref{eq:C-rade-gaussian}. 
\end{proposition} 

\begin{lemma}[Symmetrization]
	\label{lem:symmetrization} 
	Let  $q\in[1,\infty)$, $M\in\bbN$, $(r_j)_{j=1}^M$ be 
	a Rademacher family on a complete probability space 
	$(\Omegat,\cAt,\bbPt)$, 
	and let $\eta_1,\ldots,\eta_M\in L_q(\Omegat;F)$ be 
	centered random variables, 
	i.e., $\bbEt[\eta_j]=0$ for $1\leq j \leq M$, 
	taking values 
	in a real Banach space $(F,\norm{\,\cdot\,}{F})$ 
	such that $\eta_1,\ldots,\eta_M, r_1, \ldots, r_M$ are independent. 
	Then, 
	% 
	\[
	\biggl\| 
	\sum_{j=1}^M \eta_j 
	\biggr\|_{L_q(\Omegat;F)} 
	\leq 
	2 \, 
	\biggl\| 
	\sum_{j=1}^M r_j \eta_j 
	\biggr\|_{L_q(\Omegat;F)} . 
	\]
\end{lemma} 

\begin{definition}[Kahane--Khintchine constants]
	\label{def:Kqp-constant}
	Assume that $p,q\in[1,\infty)$. 
	The  $(q,p)$ Kahane--Khintchine constant $K_{q,p}$ 
	is the smallest constant $K\in(0,\infty)$ 
	such that 
	for any Rademacher family $(r_j)_{j\in\bbN}$ 
	on a complete probability space 
	$(\Omegat,\cAt,\bbPt)$, 
	for any real Banach space $(F,\norm{\,\cdot\,}{F})$,  
	for all $n\in\bbN$, 
	and every $x_1,\ldots,x_n\in F$, 
	%
	\begin{equation}\label{eq:kahane-khintchine} 
	\biggl\| 
	\sum_{j=1}^n r_j x_j
	\biggr\|_{L_q(\Omegat;F)} 
	\leq K \,
	\biggl\| 
	\sum_{j=1}^n r_j x_j
	\biggr\|_{L_p(\Omegat;F)}. 
	\end{equation} 
\end{definition} 

\begin{definition}[Type $p$ constant]\label{def:type-p-constant}
	A real Banach space $(F,\norm{\,\cdot\,}{F})$ 
	is said to be of \emph{(Rademacher) type} $p\in[1,2]$ 
	if there exists a constant $\tau\in(0,\infty)$ 
	such that 
	for any Rademacher family $(r_j)_{j\in\bbN}$ 
	on a complete probability space 
	$(\Omegat,\cAt,\bbPt)$ 
	(with expectation $\bbEt$), 
	for every $n\in\bbN$, 
	and all vectors $x_1,\ldots,x_n\in F$, 
	%
	\begin{equation}\label{eq:type} 
	\biggl\| 
	\sum_{j=1}^n r_j x_j 
	\biggr\|_{L_p(\Omegat;F)}  
	=
	\Biggl( 
	\bbEt   
	\Biggl[ 
	\biggl\|
	\sum_{j=1}^n r_j x_j 
	\biggr\|_F^p 
	\Biggr] 
	\Biggr)^{\nicefrac{1}{p}} 
	\leq 
	\tau 
	\Biggl( 
	\sum_{j=1}^n \norm{x_j}{F}^p 
	\Biggr)^{\nicefrac{1}{p}} \!. 
	\end{equation} 
	%
	In this case, the smallest constant $\tau\in(0,\infty)$ 
	in \eqref{eq:type} 
	is called the \emph{type $p$ constant of $F$} 
	and will be denoted by $\tau_p(F)$. 
\end{definition}

	\begin{equation}\label{eq:khintchine} 
	\hspace{-1.5mm} 
	A_q 
	\Biggl( \sum_{j=1}^n a_j^2 \Biggr)^{\nicefrac{1}{2}} 
	\leq 
	\Biggl( \bbEt \Biggl[ \biggl| \sum_{j=1}^n r_j a_j \biggr|^q \Biggr] \Biggr)^{\nicefrac{1}{q}} 
	= 
	\biggl\| 
	\sum_{j=1}^n r_j a_j 
	\biggr\|_{L_q(\Omegat;\bbR)}  
	\leq 
	B_q 
	\Biggl( \sum_{j=1}^n a_j^2 \Biggr)^{\nicefrac{1}{2}} 
	\hspace{-1.5mm} 
	\end{equation} 

\begin{theorem}\label{thm:MC-k} 
	Assume that  
	$(E,\norm{\,\cdot\,}{E})$ 
	is of Rademacher type $p\in[1,2]$. 
	Let $q\in[p,\infty)$, $k,M\in\bbN$ 
	and $\xi_1,\ldots,\xi_M\in L_{kq}(\Omega;E)$ be 
	independent and identically distributed 
	$E$-valued random variables. 
	Then, 
	%
	\begin{equation}\label{eq:thm:MC-k}  	
		\biggl\| 
		\Mk[\xi_1] 
		-
		\frac{1}{M} 
		\sum_{j=1}^M 
		\xk \xi_j 
		\biggr\|_{L_q(\Omega;\xkseps E)} 
		\leq 
		\CSL \,  
		M^{-\left( 1-\frac{1}{p} \right)} 
		\norm{ \xi_1 }{L_{kq}(\Omega; E)}^k .  
	\end{equation} 
	%
	Here, we recall  
	the constant $B_q\in(0,\infty)$   
	from the classical 
	Khintchine inequalities \eqref{eq:khintchine},  
	as well as the Kahane--Khintchine constant $K_{q,p}$   
	and type $p$ constant $\tau_p(E)$  
	from Definitions~\ref{def:Kqp-constant} and~\ref{def:type-p-constant}, 
	respectively, and set 
	%
	\begin{equation}\label{eq:CSL} 
	\CSL := 2 (2 k  K_{q,p} \tau_p(E) + B_q). 
	\end{equation}
\end{theorem} 

\begin{proof} 
	Assume that $(r_j)_{j=1}^M$ is 
	a Rademacher family on a complete probability space 
	$(\Omegat,\cAt,\bbPt)$ 
	and, for $j\in\{1,\ldots,M\}$, 
	let $\xivec_j\from\Omega\times\Omegat\to E$ 
	and $\rvec_j\from\Omega\times\Omegat\to \{-1,1\}$ denote the mappings 
	that satisfy $\xivec_j(\omega,\omegat) = \xi_j(\omega)$ 
	and $\rvec_j(\omega,\omegat) = r_j(\omegat)$ 
	for all $(\omega,\omegat)\in\Omega\times\Omegat$. 
	Notice that 
	on $(\Omega \times \Omegat,\cA \otimes \cAt,\bbP\otimes\bbPt)$ 
	the random variables 
	$(\rvec_j)_{j=1}^M$ form a Rademacher family and 
	$\xivec_1,\ldots,\xivec_M,\rvec_1,\ldots,\rvec_M$ 
	are independent. 
	Moreover,  
	$\xi_j \in L_{kq}(\Omega;E)$ implies that 
	$\Mk[\xi_1] \in \xkseps E$ and 
	$\xk \xivec_j \in L_q(\Omega\times\widetilde{\Omega};\xkseps E)$ 
	are well-defined. 
	We further note that by the identical distribution 
	of $\xi_1,\ldots,\xi_M$, 
	%
	\[
	(\bbE\otimes\bbEt)\bigl[  \xk \xivec_j  -  \Mk[\xi_1] \bigr]  
	= 
	\bbE\bigl[ \xk \xi_j -  \Mk[\xi_1] \bigr]  
	=
	\bbE\bigl[ \xk \xi_j \bigr] -  \Mk[\xi_j]  
	= 
	0. 
	\]  
	% 
	This shows that the independent random variables 
	$\xk \xivec_j  -  \Mk[\xi_1]\from \Omega\times\Omegat \to \xkseps E$, 
	$1\leq j\leq M$, 
	are centered. 
	Therefore, 
	we can use Lemma~\ref{lem:symmetrization}  
	on the 
	probability space 
	$(\Omega \times \Omegat,\cA \otimes \cAt,\bbP\otimes\bbPt)$ 
	and for the Banach space $\xkseps E$ to deduce that 
	% 
	\[
	\biggl\| 
	\sum_{j=1}^M 
	\bigl( \xk \xivec_j  
	-  \Mk[\xi_1] \bigr)
	\biggr\|_{L_q(\Omega\times\widetilde{\Omega};\xkseps E)} 
	\leq 
	2\, 
	\biggl\| 
	\sum_{j=1}^M 
	\rvec_j 
	\bigl( \xk \xivec_j  
	-  \Mk[\xi_1] \bigr)
	\biggr\|_{L_q(\Omega\times\widetilde{\Omega};\xkseps E)} . 
	\]
	% 
	By the triangle inequality 
	on $L_q(\Omega\times\widetilde{\Omega};\xkseps E)$ 
	we then obtain 
	that 
	% 
	\begin{equation}\label{eq:proof:thm:MC-k:A-B} 
	\begin{split}  
	\biggl\| 
	\sum_{j=1}^M 
	&\bigl(\xk \xi_j 
	-  \Mk[\xi_1] \bigr)
	\biggr\|_{L_q(\Omega;\xkseps E)} 
	= 
	\biggl\| 
	\sum_{j=1}^M 
	\bigl( \xk \xivec_j  
	-  \Mk[\xi_1] \bigr)
	\biggr\|_{L_q(\Omega\times\widetilde{\Omega};\xkseps E)} 
	\\
	&\leq 
	2\, 
	\biggl\| 
	\sum_{j=1}^M 
	\rvec_j 
	\bigl( \xk \xivec_j  
	-  \Mk[\xi_1] \bigr)
	\biggr\|_{L_q(\Omega\times\widetilde{\Omega};\xkseps E)}  
	\\
	&\leq 
	2\, 
	\biggl\| 
	\sum_{j=1}^M 
	\rvec_j 
	\xk\! \xivec_j  
	\biggr\|_{L_q(\Omega\times\widetilde{\Omega};\xkseps E)}   
	+ 
	2\, 
	\biggl\| 
	\sum_{j=1}^M 
	\rvec_j 
	\Mk[\xi_1] 
	\biggr\|_{L_q(\Omega\times\widetilde{\Omega};\xkseps E)}   
	\\
	&=: 
	2\text{(A)} + 2\text{(B)}. 
	\end{split} 
	\end{equation} 
	%
	
	To bound term (A) from above, 
	we apply Fubini's theorem and 
	the Kahane--Khintchine inequality~\eqref{eq:kahane-khintchine} 
	for the Banach space $F:=\xkseps E$ 
	and obtain that 
	%
	\begin{align*} 
	\text{(A)} 
	&=
	\Biggl( 
	\int_\Omega 
	\biggl\| 
	\sum_{j=1}^M 
	r_j (\,\cdot\,) 
	\xk\! \xi_j(\omega)  
	\biggr\|_{L_q(\Omegat;\xkseps E)}^q  
	\rd\bbP(\omega)
	\Biggr)^{\nicefrac{1}{q}} 
	\\
	&\leq 
	K_{q,p} 
	\Biggl( 
	\int_\Omega 
	\biggl\| 
	\sum_{j=1}^M 
	r_j (\,\cdot\,) 
	\xk\! \xi_j(\omega) 
	\biggr\|_{L_p(\Omegat;\xkseps E)}^q  
	\rd\bbP(\omega)
	\Biggr)^{\nicefrac{1}{q}} \!. 
	\end{align*} 
	% 
	Upon inserting the definition~\eqref{eq:sepsnorm}  
	of the symmetric injective tensor norm,  
	we use Proposition~\ref{prop:rade-gaussian}\ref{prop:rade-gaussian-poly}
	for the convex increasing function $G(t) := t^p$, $t\geq0$,
	and 
	the fact that the Banach space $E$ has type $p\in[1,2]$ 
	with type constant $\tau_p(E)\in(0,\infty)$, 
	to conclude that 
	% 
	{\allowdisplaybreaks
	\begin{align*} 
	\text{(A)} 
	&\leq 
	K_{q,p} 
	\Biggl( 
	\int_\Omega 
	\Biggl( 
	\bbEt 
	\Biggl[ 
	\biggl( 
	\sup_{f\in B_{\dual{E}}} 
	\biggl| 
	\sum_{j=1}^M 
	r_j (\,\cdot\,) 
	f\bigl( \xi_j(\omega) \bigr)^k 
	\biggr| \biggr)^p 
	\Biggr] 
	\Biggr)^{\nicefrac{q}{p}} 
	\rd\bbP(\omega)
	\Biggr)^{\nicefrac{1}{q}}  
	\\
	&\leq  
	2 k K_{q,p}  
	\Biggl( 
	\int_\Omega 
	\Biggl( 
	\bbEt 
	\Biggl[ 
	\biggl\| 
	\sum_{j=1}^M 
	r_j (\,\cdot\,) 
	\norm{ \xi_j(\omega) }{E}^{k-1} 
	\xi_j(\omega)
	\biggr\|_E^p 
	\Biggr] 
	\Biggr)^{\nicefrac{q}{p}} 
	\rd\bbP(\omega)
	\Biggr)^{\nicefrac{1}{q}} 
	\\
	\text{(A)} 
	&\leq 
	2 k K_{q,p} \tau_p(E) 
	\Biggl( 
	\int_\Omega 
	\Biggl( 
	\sum_{j=1}^M 
	\norm{ \xi_j(\omega) }{E}^{kp} 
	\Biggr)^{\nicefrac{q}{p}} 
	\rd\bbP(\omega)
	\Biggr)^{\nicefrac{1}{q}} 
	\\
	&= 
	2 k K_{q,p} \tau_p(E) 
	\biggl\| 
	\sum_{j=1}^M 
	\norm{\xi_j}{E}^{kp} 
	\biggr\|_{L_{\nicefrac{q}{p}}(\Omega;\bbR)}^{\nicefrac{1}{p}}. 
	\end{align*} 
	% 
	Since} $q\geq p$, 
	we can use the triangle inequality 
	on $L_{\nicefrac{q}{p}}(\Omega;\bbR)$, yielding 
	% 
	\begin{equation}\label{eq:proof:thm:MC-k:A} 
	\text{(A)} 
	\leq 
	2 k K_{q,p} \tau_p(E) 
	\Biggl(  
	\sum_{j=1}^M 
	\norm{\xi_j}{L_{kq}(\Omega;E)}^{kp} 
	\biggr)^{\nicefrac{1}{p}} 
	=  
	2 k K_{q,p} \tau_p(E) 
	M^{\nicefrac{1}{p}} 
	\norm{\xi_1}{L_{kq}(\Omega; E)}^k, 
	\end{equation} 
	% 
	where we also used the  
	identical distribution of $\xi_1,\ldots,\xi_M$. 
	
	For term (B) we use the estimate  
	% 
	\[
	\bigl\| \Mk[\xi_1] \bigr\|_{\eps_s}  
	= 
	\bigl\| \bbE\bigl[ \xk \xi_1 \bigr] \bigr\|_{\eps_s} 
	\leq 
	\bbE\bigl[ \norm{ {\xk} \xi_1 }{\eps_s} \bigr]
	= 
	\bbE\bigl[ \norm{ \xi_1 }{E}^k \bigr] 
	\leq 
	\norm{ \xi_1 }{L_{kq}(\Omega;E)}^k , 
	\]
	% 
	as well as the classical 
	Khintchine inequalities \eqref{eq:khintchine}
	so that 
	% 
	\[
	\biggl\| 
	\sum_{j=1}^M 
	r_j 
	\biggr\|_{L_q(\widetilde{\Omega};\bbR)} 
	\leq 
	B_q \, M^{\nicefrac{1}{2}} 
	\leq 
	B_q \, M^{\nicefrac{1}{p}} \!, 
	\]
	%
	and conclude that 
	% 
	\begin{equation}\label{eq:proof:thm:MC-k:B}   
	\begin{split} 
	\text{(B)}  
	&= 
	\biggl\| 
	\Mk[\xi_1] 
	\sum_{j=1}^M 
	\rvec_j 
	\biggr\|_{L_q(\Omega\times\widetilde{\Omega};\xkseps E)}   
	= 
	\Biggl( 
	\int_{\widetilde{\Omega}} 
	\biggl\| 
	\Mk[\xi_1] 
	\sum_{j=1}^M 
	r_j (\omegat) 
	\biggr\|_{\eps_s}^q   
	\, \rd\bbPt(\omegat) 
	\Biggr)^{\nicefrac{1}{q}} 
	\\
	&=
	\bigl\| 
	\Mk[\xi_1]  
	\bigr\|_{\eps_s} 
	\Biggl( 
	\int_{\widetilde{\Omega}} 
	\biggl| 
	\sum_{j=1}^M 
	r_j(\widetilde{\omega}) 
	\biggr|^q 
	\, \rd\bbPt(\omegat) 
	\Biggr)^{\nicefrac{1}{q}} 
	\leq 
	B_q \, M^{\nicefrac{1}{p}} 
	\norm{ \xi_1 }{L_{kq}(\Omega; E)}^k . 
	\end{split} 
	\end{equation} 
	%
	Finally, combining \eqref{eq:proof:thm:MC-k:A-B}, 
	\eqref{eq:proof:thm:MC-k:A} and 
	\eqref{eq:proof:thm:MC-k:B} shows that 
	% 
	\begin{align*} 
	\biggl\| 
	\Mk[\xi_1] 
	-
	\frac{1}{M} 
	\sum_{j=1}^M 
	\xk \xi_j 
	\biggr\|_{L_q(\Omega;\xkseps E)} 
	&=
	M^{-1} \, 
	\biggl\| 
	\sum_{j=1}^M 
	\bigl(\xk \xi_j -  \Mk[\xi_1] \bigr)
	\biggr\|_{L_q(\Omega;\xkseps E)} 
	\\
	&\leq 
	2 
	( 2k K_{q,p} \tau_p(E) + B_q) \, 
	M^{-\left(1-\frac{1}{p}\right)} 
	\norm{ \xi_1 }{L_{kq}(\Omega;E)}^k, 
	\end{align*} 
	%
	which along with 
	the definition \eqref{eq:CSL} 
	of~$\CSL$ 
	completes the proof 
	of \eqref{eq:thm:MC-k}. 
\end{proof} 

\end{document}
